\chapter{The Mismatch String Kernel}\label{ch:kernel}

\section{Kernel Methods: A Brief Introduction}\label{sec:kernel-intro}

% Definition of a kernel function: $K(x, x') = \langle \Phi(x), \Phi(x') \rangle_{\mathcal{H}}$
% The kernel trick: computing inner products in a high-dimensional feature space
% without explicitly constructing the feature map $\Phi$.
% Mercer's theorem and the characterisation of valid (positive semi-definite) kernels.
% Why kernels are well-suited for structured data (strings, graphs, sets).

\section{String Kernels}\label{sec:string-kernels}

\subsection{The Spectrum Kernel}\label{sec:spectrum}
% The $k$-spectrum kernel (Leslie et al., 2002):
% $\Phi_k(x) = (\phi_\alpha(x))_{\alpha \in \Sigma^k}$, where $\phi_\alpha(x)$
% counts the number of occurrences of $k$-mer $\alpha$ in string $x$.
% Limitations: exact matching only, sensitive to single-nucleotide variations.

\subsection{The Mismatch Kernel}\label{sec:mismatch-kernel}
% The $(k, m)$-mismatch kernel (Leslie et al., 2004):
% For each observed $k$-mer, increment the count of all $k$-mers within Hamming
% distance $m$ from it. This neighbourhood expansion makes the kernel robust
% to substitutions.
%
% Formal definition:
% $N(k, m)(\alpha) = \{ \beta \in \Sigma^k : d_H(\alpha, \beta) \leq m \}$
% $\Phi_{k,m}(x) = \sum_{i=1}^{|x|-k+1} \mathbf{1}_{N(k,m)(x[i:i+k])}$
%
% Alphabet: the extended 6-letter epigenetic alphabet $\Sigma = \{A, C, G, T, M, H\}$

\subsection{Weighted Mismatch Neighbourhood}\label{sec:weighted-mismatch}
% Extension used in this work: instead of binary neighbourhood membership,
% each neighbour at Hamming distance $d$ contributes weight $\lambda^d$
% where $\lambda \in [0, 1]$ is a decay factor (mismatch\_decay).
% This gives higher importance to exact matches and near-matches.

\section{Multiple Kernel Learning}\label{sec:mkl}
% Combining kernels across multiple values of $k$:
% $K_{\text{MKL}}(x, x') = \sum_{k=k_{\min}}^{k_{\max}} w_k \, K_k(x, x')$
% Weight generation strategy: linear ramp with noise suppression for short $k$-mers.

\section{Gram Matrix Computation}\label{sec:gram}

\subsection{Symmetric Gram Matrix (Training)}\label{sec:gram-symmetric}
% $K_{ij} = \langle \Phi(x_i), \Phi(x_j) \rangle$ for all pairs in the training set.
% Exploiting symmetry: only computing the upper triangle.
% Kernel normalisation: $\hat{K}_{ij} = K_{ij} / \sqrt{K_{ii} K_{jj}}$

\subsection{Asymmetric Gram Matrix (Inference)}\label{sec:gram-asymmetric}
% At test time we need $K_{\text{test}}[i, j] = \langle \Phi(x_i^{\text{test}}), \Phi(x_j^{\text{train}}) \rangle$.
% Proper normalisation requires the test self-norms $K^{\text{test}}_{ii}$, which
% differ from training norms. Describe the vocabulary alignment issue and the
% full-vocabulary optimisation.

\section{Algorithmic Optimisations}\label{sec:optimisations}

\subsection{Sparse Feature Representation}\label{sec:sparse}
% For large $k$ and small $m$, the feature vectors are extremely sparse.
% Using scipy.sparse CSR matrices reduces memory and enables fast dot products.

\subsection{LRU-Cached Neighbourhood Generation}\label{sec:lru-cache}
% The function \texttt{generate\_mismatch\_neighborhood} is decorated with
% \texttt{@lru\_cache}: the same $k$-mer neighbourhood is never recomputed.

\subsection{Pre-Enumerated Vocabulary}\label{sec:vocab}
% \texttt{build\_full\_vocabulary} generates all $|\Sigma|^k$ possible $k$-mers
% deterministically, allowing parallel chunks to operate with no shared mutable
% state and guaranteeing consistent column indices.

\subsection{Block-Partitioned Matrix Multiplication}\label{sec:block-parallel}
% For large $N$, the Gram matrix is tiled into blocks of size $B \times B$.
% Each block $K[r_s:r_e, c_s:c_e] = X[r_s:r_e] \cdot X[c_s:c_e]^T$ is
% computed independently via joblib, avoiding materialisation of a single
% dense $N \times N$ product.

\subsection{BLAS Thread Pinning}\label{sec:blas}
% Setting \texttt{OMP\_NUM\_THREADS=1} and equivalent variables before parallel
% kernel computation prevents hidden thread contention between joblib workers
% and the underlying BLAS implementation.
